% =============================================================================
% SECTION 5A: VERIFICATION, VALIDATION, AND UNCERTAINTY QUANTIFICATION
% =============================================================================

\section{Verification, Validation, and Uncertainty Quantification}
\label{sec:vvuq}

The structural mechanics layer introduces continuum-based models and probabilistic reasoning that require rigorous verification and validation to ensure correctness and physical fidelity.

\subsection{Verification: Numerical Correctness}

Verification ensures that the implemented algorithms correctly solve the mathematical models.

\textbf{Beam Solver Convergence.}
The Euler--Bernoulli beam solver is verified against closed-form analytical solutions for canonical loading cases.
For a simply-supported beam with uniform load $w$, the maximum deflection is:
\[
\delta_{\max} = \frac{5wL^4}{384EI}
\]
The numerical solver (finite-difference or finite-element discretization) must achieve relative error $< 0.1\%$ compared to this closed form.
Convergence tests with mesh refinement ensure spatial discretization errors decay at expected rates.

\textbf{Load Conservation.}
Tributary area assignments must conserve total roof and floor area:
\[
\sum_{i=1}^{N} A_{\text{trib},i} = A_{\text{total}}
\]
where $A_{\text{trib},i}$ is the tributary area assigned to header/beam $i$.
This conservation law is verified programmatically for every generated hypothesis.

\textbf{Uncertainty Propagation Convergence.}
Monte Carlo sampling converges as $N \to \infty$.
For probability-of-failure estimates, convergence is tested by comparing $N = 1000$ vs.\ $N = 10000$ samples.
Relative error in $P(\text{failure})$ must be $< 5\%$ to ensure statistical stability.

\subsection{Validation: Physical Correctness}

Validation ensures that the models and inferences correspond to real-world structural behavior.

\textbf{IRC Span Table Comparison.}
Generated header assemblies are compared against prescriptive IRC span tables~\cite{irc2021}.
For standard load scenarios (40 psf roof + 10 psf dead load), the system must select headers that match or exceed IRC prescriptions for at least 90\% of cases.
Discrepancies are flagged for manual review or hypothesis rejection.

\textbf{Engineered Design Comparison.}
A validation dataset of 20--30 residential plan sets with stamped structural designs is used to test hypothesis generation.
Success criterion: the true engineered design (or a mechanically equivalent alternative) appears in the top 3 ranked hypotheses for at least 80\% of cases.
This metric validates both the hypothesis generation strategy and the ranking heuristics.

\textbf{Ground Truth Annotations.}
For each validation plan, expert estimators provide:
(1) identified load-bearing walls,
(2) header sizing rationale,
(3) tributary area sketches.
These annotations serve as oracle data for evaluating inference accuracy.

\subsection{Uncertainty Quantification Framework}

Structural loads and material properties are inherently uncertain. The system models these uncertainties explicitly.

\textbf{Random Variable Specification.}
\begin{itemize}
  \item \textbf{Tributary width:} $W \sim \text{Uniform}(W_{\min}, W_{\max})$, where bounds reflect geometric ambiguity.
  \item \textbf{Roof load:} $q_{\text{roof}} \sim \mathcal{N}(\mu_{\text{roof}}, \sigma_{\text{roof}})$, with $\mu$ from code snow/wind requirements and $\sigma$ reflecting uncertainty.
  \item \textbf{Floor load:} $q_{\text{floor}} \sim \mathcal{N}(\mu_{\text{floor}}, \sigma_{\text{floor}})$.
  \item \textbf{Material stiffness:} $E \sim \mathcal{N}(\mu_E, \sigma_E)$, accounting for wood grade variability.
\end{itemize}

\textbf{Monte Carlo Propagation.}
For each structural hypothesis $h$ and candidate header assembly $a$:
\begin{enumerate}
  \item Sample $N = 1000$ realizations: $(W^{(j)}, q_{\text{roof}}^{(j)}, q_{\text{floor}}^{(j)}, E^{(j)})$ for $j = 1, \ldots, N$.
  \item For each realization, solve the beam PDE to compute $(\sigma^{(j)}, \tau^{(j)}, \delta^{(j)})$.
  \item Estimate probability of failure:
  \[
  P_{\text{fail}}(a, h) \approx \frac{1}{N} \sum_{j=1}^{N} \mathbb{1}\left(\sigma^{(j)} > F_b \cup \tau^{(j)} > F_v \cup \delta^{(j)} > L/240\right)
  \]
  where $\mathbb{1}(\cdot)$ is the indicator function.
\end{enumerate}

\textbf{Robust Feasibility.}
Assemblies are accepted only if $P_{\text{fail}} < 0.01$ (1\% failure probability threshold).
This ensures that structural decisions remain feasible under realistic load and material variability.

\subsection{Hypothesis Scoring and Selection}

Multiple structural hypotheses are ranked using a multi-criteria score:
\[
S(h) = w_1 \cdot (1 - C_{\text{norm}}(h)) + w_2 \cdot (1 - P_{\text{fail,max}}(h)) + w_3 \cdot F(h)
\]
where:
\begin{itemize}
  \item $C_{\text{norm}}(h)$ is normalized material cost proxy (lumber volume + engineered member usage),
  \item $P_{\text{fail,max}}(h) = \max_{a \in h} P_{\text{fail}}(a, h)$ is the worst-case failure probability across all members,
  \item $F(h)$ is design flexibility score (fraction of interior walls that remain non-bearing),
  \item $w_1, w_2, w_3$ are user-tunable weights (default: $0.4, 0.4, 0.2$).
\end{itemize}

The top 2--3 Pareto-optimal hypotheses are presented to the user, along with tradeoff explanations (e.g., ``Hypothesis 1: lowest cost but requires interior post; Hypothesis 2: no interior posts but 15\% higher material cost'').

\subsection{Scientific Computing Contribution}

This work contributes to scientific computing in construction by:
\begin{itemize}
  \item \textbf{Automatic PDE generation:} Translating discrete geometric plans into continuous beam boundary value problems.
  \item \textbf{Reduced-order continuum models:} Using Euler--Bernoulli beam theory (a reduced-order PDE derived from 3D elasticity) for real-time structural feasibility.
  \item \textbf{Hybrid discrete-continuous optimization:} Combining graph-based load path inference with PDE-constrained optimization.
  \item \textbf{Uncertainty-aware decision-making:} Integrating Monte Carlo UQ with construction workflows for robust material selection.
\end{itemize}

These contributions bridge the gap between theoretical structural mechanics and practical construction planning, enabling physics-informed automation without requiring full structural engineering analysis.
